\documentclass[a4paper]{arrowhead}

\usepackage[yyyymmdd]{datetime}
\usepackage{etoolbox}
\usepackage[utf8]{inputenc}
\usepackage[T1]{fontenc}
\usepackage{multirow}
\usepackage{longtable}
\usepackage{tikz}
\usepackage{adjustbox}
\AtBeginEnvironment{longtable}{\small}
\BeforeBeginEnvironment{tikzpicture}{\begin{adjustbox}{max width=\textwidth}}
\AfterEndEnvironment{tikzpicture}{\end{adjustbox}}
\usetikzlibrary{positioning,shapes.multipart,arrows.meta}

\renewcommand{\dateseparator}{-}
\setlength{\emergencystretch}{3em}
\lstset{language=HTTP,breakatwhitespace=false,columns=fullflexible}

\newcommand{\fparam}[1]{\textit{\textcolor{ArrowheadBlue}{#1}}}

%% Special references
\newcommand{\mref}[2][]{{\textcolor{ArrowheadPurple}{\hyperref[sec:model:#2]{\ifstrempty{#1}{#2}{#1}}}}}
\newcommand{\oref}[1]{{\textcolor{ArrowheadBlue}{\hyperref[sec:operations:#1]{#1}}}}
\newcommand{\pdef}[1]{{\textcolor{ArrowheadGrey}{#1 \label{sec:model:primitives:#1} \label{sec:model:primitives:#1s}}}}
\newcommand{\pref}[1]{{\textcolor{ArrowheadGrey}{\hyperref[sec:model:primitives:#1]{#1}}}}

% as template_idd.tex; the input and output arguments are passed already formatted (\mref or \pref)
\newrobustcmd\osubsection[5]{
  \addtocounter{subsection}{1}
  \addcontentsline{toc}{subsection}{\protect\numberline{\thesubsection}operation \textcolor{ArrowheadBlue}{#1}}
  \subsection*{
    \thesubsection\quad
    #2 \textcolor{ArrowheadPurple}{#3} \\
    \small
    \hspace*{0.075\textwidth}\begin{minipage}{0.1\textwidth}
      \vspace*{1mm}
      Operation: \\
      \notblank{#4}{Input: \\}{}
      \notblank{#5}{Output: \\}{}
    \end{minipage}%
    \begin{minipage}{0.825\textwidth}
      \vspace*{1mm}
      \textcolor{ArrowheadBlue}{#1} \\
      \notblank{#4}{#4 \\}{}
      \notblank{#5}{#5 \\}{}
    \end{minipage}
  }
  \label{sec:operations:#1}
}
\newrobustcmd\msubsection[2]{
  \addtocounter{subsection}{1}
  \addcontentsline{toc}{subsection}{\protect\numberline{\thesubsection}#1 \textcolor{ArrowheadPurple}{#2}}
  \subsection*{\thesubsection\quad#1 \textcolor{ArrowheadPurple}{#2}}
  \label{sec:model:#2}
}
%%

\begin{document}
\sloppy

%% Arrowhead Document Properties
\ArrowheadTitle{authorizationTokenManagement HTTP/TLS/JSON}
\ArrowheadServiceID{authorizationTokenManagement}
\ArrowheadType{Interface Design Description}
\ArrowheadTypeShort{IDD}
\ArrowheadVersion{5.2.2}
\ArrowheadDate{2026-10-22}
\ArrowheadAuthor{Jerker Delsing}
\ArrowheadStatus{RELEASE CANDIDATE}
\ArrowheadContact{jerker.delsing@ltu.se}
\ArrowheadFooter{\href{www.arrowhead.eu}{www.arrowhead.eu}}
\ArrowheadSetup
%%

%% Front Page
\begin{center}
  \vspace*{1cm}
  \huge{\arrowtitle}

  \vspace*{0.2cm}
  \LARGE{\arrowtype}
  \vspace*{1cm}

  \Large{Service ID: \textit{"\arrowid"}}

  \vspace*{\fill}

  % Front Page Image

  \vspace*{1cm}
  \vspace*{\fill}

  % Front Page Abstract
  \begin{abstract}
    This document describes the HTTP/1.1 protocol with TLS payload security and JSON payload encoding variant of the \texttt{authorization\allowbreak{}Token\allowbreak{}Management} service of the ConsumerAuthorization system.
  \end{abstract}

  \vspace*{1cm}
\end{center}

\newpage
%%

%% Table of Contents
\tableofcontents
\newpage
%%

\section{Overview}
\label{sec:overview}

This document describes the \texttt{authorization\allowbreak{}Token\allowbreak{}Management} service interface of the ConsumerAuthorization system, which provides operator management of authorization tokens and of provider encryption keys.

It's implemented using protocol, encoding as stated in below table:
\begin{table}[ht!]
  \centering
  \begin{tabular}{|l|l|l|}
    \hline
    \rowcolor{gray!33} Profile type & Type & Version \\ \hline
    Transfer protocol & HTTP & 1.1 \cite{rfc9112} \\ \hline
    Data encryption & TLS & 1.3 \cite{rfc8446}, 1.2 \cite{rfc5246} \\ \hline
    Encoding & JSON & RFC 8259 \cite{rfc8259} \\ \hline
    Compression & N/A & - \\ \hline
    Semantics & N/A & - \\ \hline
    Ontology & N/A & - \\ \hline
  \end{tabular}
  \caption{Communication and semantics details used for the \texttt{authorization\allowbreak{}Token\allowbreak{}Management} service interface}
  \label{tab:comunication_semantics_profile}
\end{table}

This document provides the Interface Design Description IDD to the \textit{authorizationTokenManagement -- Service Description} document \cite{ca-authorizationTokenManagement-sd}.
For further details about how this service is meant to be used, please consult that document.

The rest of this document describes how to realize the \texttt{authorization\allowbreak{}Token\allowbreak{}Management} service
interface in detail. Both in terms of its operations (Section \ref{sec:operations}) and its information model (Section \ref{sec:model}).

\subsection{Interface registration}

The service is registered in the ServiceRegistry with the interface template \texttt{generic\_\allowbreak{}https} (\texttt{generic\_\allowbreak{}http} when TLS is disabled), access port 8445 and base path \texttt{/\allowbreak{}consumerauthorization/\allowbreak{}authorization/\allowbreak{}mgmt/\allowbreak{}token}. The \texttt{operations} property of the interface lists the operations below with their method and path relative to the base path (path parameters are placeholders such as \texttt{\{id\}} that the consumer fills in):

\noindent\begin{longtable}{| p{5.2cm} | p{1.6cm} | p{8.4cm} |} \hline
\rowcolor{gray!33} Operation & Method & Path \\ \hline \endhead
\texttt{generate-\allowbreak{}tokens} & POST & \texttt{/\allowbreak{}generate} \\ \hline
\texttt{query-\allowbreak{}tokens} & POST & \texttt{/\allowbreak{}query} \\ \hline
\texttt{revoke-\allowbreak{}tokens} & DELETE & \texttt{/\allowbreak{}revoke} \\ \hline
\texttt{add-\allowbreak{}encryption-\allowbreak{}keys} & POST & \texttt{/\allowbreak{}encryption-\allowbreak{}key} \\ \hline
\texttt{remove-\allowbreak{}encryption-\allowbreak{}keys} & DELETE & \texttt{/\allowbreak{}encryption-\allowbreak{}key} \\ \hline
\end{longtable}

\newpage

\section{Service Operations}
\label{sec:operations}

The interfaces of the \texttt{authorization\allowbreak{}Token\allowbreak{}Management} service, its operations, data models and implementation are provided below.
In particular, each subsection first names the HTTP method and path used to call the operation, after which it names the abstract operation from the \texttt{authorization\allowbreak{}Token\allowbreak{}Management} SD document, as well as input and output types.

All operations require the requester's identity, presented according to the provider's authentication policy: the client certificate (mutual TLS) and, with the \texttt{outsourced} policy, \texttt{Authorization: Bearer IDENTITY-TOKEN//<token>}; systems holding the Authentication secret key send \texttt{Bearer AUTHENTICATOR-KEY//<system>//<hash>}, and with the \texttt{declared} policy \texttt{Bearer SYSTEM//<system>} is sent. Request and response bodies are \texttt{application/json}.
Errors are answered with an \mref{ErrorMessage} body and the status codes \texttt{400 Bad Request} (invalid parameter), \texttt{401 Unauthorized} (missing or invalid identity), \texttt{403 Forbidden} (not permitted, e.g. by the management or blacklist filter), \texttt{404 Not Found}, \texttt{408 Request Timeout}, \texttt{423 Locked}, \texttt{500 Internal Server Error} and \texttt{503 Service Unavailable} (a system it depends on failed).

\osubsection{generate-tokens}{POST}{/\allowbreak{}consumerauthorization/\allowbreak{}authorization/\allowbreak{}mgmt/\allowbreak{}token/\allowbreak{}generate?unbound=\fparam{\{unbound\}}}{\mref[Authorization\allowbreak{}Token\allowbreak{}Generation\allowbreak{}Mgmt\allowbreak{}List\allowbreak{}Request]{AuthorizationTokenGenerationMgmtListRequest}}{\mref[Authorization\allowbreak{}Token\allowbreak{}Mgmt\allowbreak{}List\allowbreak{}Response]{AuthorizationTokenMgmtListResponse}}

Generates tokens on behalf of consumers. With \texttt{unbound=\allowbreak{}true} the authorization check is skipped (only for systems in \texttt{unbounded.\allowbreak{}token.\allowbreak{}generation.\allowbreak{}whitelist}).

Query parameters: \fparam{unbound} (boolean, optional, default \texttt{false}). List values are repeated or comma-separated.

A successful call returns \texttt{201 Created} with the output in the body.

\begin{lstlisting}[label={lst:generate-tokens},caption={An \oref{generate-tokens} invocation.}]
POST /consumerauthorization/authorization/mgmt/token/generate?unbound=<unbound> HTTP/1.1
Host: consumerauthorization:8445
Authorization: Bearer IDENTITY-TOKEN//<token>
Content-Type: application/json

{
  "list": [
    {
      "tokenVariant": "TIME_LIMITED_TOKEN_AUTH",
      "targetType": "SERVICE_DEF",
      "consumerCloud": "LOCAL",
      "consumer": "TemperatureDisplay",
      "provider": "TemperatureSensor",
      "target": "temperatureReading",
      "scope": "read",
      "expiresAt": "2026-10-22T12:00:00Z",
      "usageLimit": 5
    }
  ]
}
\end{lstlisting}

\begin{lstlisting}[label={lst:generate-tokens-resp},caption={An \oref{generate-tokens} response.}]
{
  "entries": [
    {
      "tokenType": "TIME_LIMITED_TOKEN",
      "variant": "TIME_LIMITED_TOKEN_AUTH",
      "token": "<token>",
      "tokenReference": "<token reference>",
      "requester": "DynamicServiceOrchestration",
      "consumerCloud": "LOCAL",
      "consumer": "TemperatureDisplay",
      "provider": "TemperatureSensor",
      "targetType": "SERVICE_DEF",
      "target": "temperatureReading",
      "scope": "read",
      "createdAt": "2026-10-22T12:00:00Z",
      "usageLimit": 5,
      "usageLeft": 5,
      "expiresAt": "2026-10-22T12:00:00Z"
    }
  ],
  "count": 1
}
\end{lstlisting}

\osubsection{query-tokens}{POST}{/\allowbreak{}consumerauthorization/\allowbreak{}authorization/\allowbreak{}mgmt/\allowbreak{}token/\allowbreak{}query}{\mref[Authorization\allowbreak{}Token\allowbreak{}Query\allowbreak{}Request]{AuthorizationTokenQueryRequest}}{\mref[Authorization\allowbreak{}Token\allowbreak{}Mgmt\allowbreak{}List\allowbreak{}Response]{AuthorizationTokenMgmtListResponse}}

Returns a page of tokens matching the filter.

A successful call returns \texttt{200 OK} with the output in the body.

\begin{lstlisting}[label={lst:query-tokens},caption={An \oref{query-tokens} invocation.}]
POST /consumerauthorization/authorization/mgmt/token/query HTTP/1.1
Host: consumerauthorization:8445
Authorization: Bearer IDENTITY-TOKEN//<token>
Content-Type: application/json

{
  "pagination": {
    "page": 0,
    "size": 10,
    "direction": "ASC",
    "sortField": "createdAt"
  },
  "requester": "DynamicServiceOrchestration",
  "tokenType": "TIME_LIMITED_TOKEN",
  "consumerCloud": "LOCAL",
  "consumer": "TemperatureDisplay",
  "provider": "TemperatureSensor",
  "targetType": "SERVICE_DEF",
  "target": "temperatureReading"
}
\end{lstlisting}

\begin{lstlisting}[label={lst:query-tokens-resp},caption={An \oref{query-tokens} response.}]
{
  "entries": [
    {
      "tokenType": "TIME_LIMITED_TOKEN",
      "variant": "TIME_LIMITED_TOKEN_AUTH",
      "token": "<token>",
      "tokenReference": "<token reference>",
      "requester": "DynamicServiceOrchestration",
      "consumerCloud": "LOCAL",
      "consumer": "TemperatureDisplay",
      "provider": "TemperatureSensor",
      "targetType": "SERVICE_DEF",
      "target": "temperatureReading",
      "scope": "read",
      "createdAt": "2026-10-22T12:00:00Z",
      "usageLimit": 5,
      "usageLeft": 5,
      "expiresAt": "2026-10-22T12:00:00Z"
    }
  ],
  "count": 1
}
\end{lstlisting}

\osubsection{revoke-tokens}{DELETE}{/\allowbreak{}consumerauthorization/\allowbreak{}authorization/\allowbreak{}mgmt/\allowbreak{}token/\allowbreak{}revoke?tokenReferences=\fparam{\{tokenReferences\}}}{\pref{Array}$<$\allowbreak{}\pref{String}$>$}{}

Revokes tokens by token reference.

Query parameters: \fparam{tokenReferences} (List\textless{}String\textgreater{}, required). List values are repeated or comma-separated.

A successful call returns \texttt{200 OK}.

\begin{lstlisting}[label={lst:revoke-tokens},caption={An \oref{revoke-tokens} invocation.}]
DELETE /consumerauthorization/authorization/mgmt/token/revoke?tokenReferences=<tokenReferences> HTTP/1.1
Host: consumerauthorization:8445
Authorization: Bearer IDENTITY-TOKEN//<token>
\end{lstlisting}

\osubsection{add-encryption-keys}{POST}{/\allowbreak{}consumerauthorization/\allowbreak{}authorization/\allowbreak{}mgmt/\allowbreak{}token/\allowbreak{}encryption-key}{\mref[Authorization\allowbreak{}Mgmt\allowbreak{}Encryption\allowbreak{}Key\allowbreak{}Registration\allowbreak{}List\allowbreak{}Request]{AuthorizationMgmtEncryptionKeyRegistrationListRequest}}{\mref[Authorization\allowbreak{}Mgmt\allowbreak{}Encryption\allowbreak{}Key\allowbreak{}List\allowbreak{}Response]{AuthorizationMgmtEncryptionKeyListResponse}}

Registers encryption keys on behalf of providers.

A successful call returns \texttt{201 Created} with the output in the body.

\begin{lstlisting}[label={lst:add-encryption-keys},caption={An \oref{add-encryption-keys} invocation.}]
POST /consumerauthorization/authorization/mgmt/token/encryption-key HTTP/1.1
Host: consumerauthorization:8445
Authorization: Bearer IDENTITY-TOKEN//<token>
Content-Type: application/json

{
  "list": [
    {
      "systemName": "TemperatureSensor",
      "key": "0123456789abcdef",
      "algorithm": "AES/ECB/PKCS5Padding"
    }
  ]
}
\end{lstlisting}

\begin{lstlisting}[label={lst:add-encryption-keys-resp},caption={An \oref{add-encryption-keys} response.}]
{
  "entries": [
    {
      "systemName": "TemperatureSensor",
      "rawKey": "0123456789abcdef",
      "algorithm": "AES/ECB/PKCS5Padding",
      "keyAdditive": "",
      "createdAt": "2026-10-22T12:00:00Z"
    }
  ],
  "count": 1
}
\end{lstlisting}

\osubsection{remove-encryption-keys}{DELETE}{/\allowbreak{}consumerauthorization/\allowbreak{}authorization/\allowbreak{}mgmt/\allowbreak{}token/\allowbreak{}encryption-key?systemNames=\fparam{\{systemNames\}}}{\pref{Array}$<$\allowbreak{}\pref{String}$>$}{}

Removes the encryption keys of the named providers.

Query parameters: \fparam{systemNames} (List\textless{}String\textgreater{}, required). List values are repeated or comma-separated.

A successful call returns \texttt{200 OK}.

\begin{lstlisting}[label={lst:remove-encryption-keys},caption={An \oref{remove-encryption-keys} invocation.}]
DELETE /consumerauthorization/authorization/mgmt/token/encryption-key?systemNames=<systemNames> HTTP/1.1
Host: consumerauthorization:8445
Authorization: Bearer IDENTITY-TOKEN//<token>
\end{lstlisting}

\newpage

\section{Data Models}
\label{sec:model}

Here, all data objects that can be part of the service calls associated with this service are listed in alphabetic order.
Note that each subsection, which describes one type of object, begins with the \textit{struct} keyword, which is meant to denote a JSON \pref{Object} that must contain certain fields, or names, with values conforming to explicitly named types.
Fields that are not needed may be omitted or \texttt{null} in requests; empty fields are omitted in responses.
As a complement to the primary types defined in this section, there is also a list of secondary types in Section \ref{sec:model:primitives}, which are used to represent things like hashes, identifiers and texts.

\msubsection{struct}{AuthorizationMgmtEncryptionKeyListResponse}

List of registered encryption keys.


\noindent\begin{tabularx}{\textwidth}{| >{\raggedright\arraybackslash}p{3.9cm} | >{\raggedright\arraybackslash\small}p{5.3cm} | X |} \hline
\rowcolor{gray!33} Object Field & Value Type & Description \\ \hline
"entries" & \pref{Array}$<$\allowbreak{}\mref[Authorization\allowbreak{}Mgmt\allowbreak{}Encryption\allowbreak{}Key\allowbreak{}Response]{AuthorizationMgmtEncryptionKeyResponse}$>$ & Matching records of the current page. \\ \hline
"count" & \pref{Number} & Total number of matching records (all pages). \\ \hline
\end{tabularx}



\msubsection{struct}{AuthorizationMgmtEncryptionKeyRegistrationListRequest}

Encryption keys to register on behalf of providers.


\noindent\begin{tabularx}{\textwidth}{| >{\raggedright\arraybackslash}p{3.9cm} | >{\raggedright\arraybackslash\small}p{5.3cm} | X |} \hline
\rowcolor{gray!33} Object Field & Value Type & Description \\ \hline
"list" & \pref{Array}$<$\allowbreak{}\mref[Authorization\allowbreak{}Mgmt\allowbreak{}Encryption\allowbreak{}Key\allowbreak{}Registration\allowbreak{}Request]{AuthorizationMgmtEncryptionKeyRegistrationRequest}$>$ & The list of request items. \\ \hline
\end{tabularx}



\msubsection{struct}{AuthorizationMgmtEncryptionKeyRegistrationRequest}

Encryption key of one provider.


\noindent\begin{tabularx}{\textwidth}{| >{\raggedright\arraybackslash}p{3.9cm} | >{\raggedright\arraybackslash\small}p{5.3cm} | X |} \hline
\rowcolor{gray!33} Object Field & Value Type & Description \\ \hline
"system\allowbreak{}Name" & \pref{String} & Provider system the key belongs to. \\ \hline
"key" & \pref{String} & The encryption key, as a string. \\ \hline
"algorithm" & \pref{String} & Encryption algorithm of the key, e.g. \texttt{AES/\allowbreak{}ECB/\allowbreak{}PKCS5Padding} or \texttt{AES/\allowbreak{}CBC/\allowbreak{}PKCS5Padding}. \\ \hline
\end{tabularx}



\msubsection{struct}{AuthorizationMgmtEncryptionKeyResponse}

A registered encryption key.


\noindent\begin{tabularx}{\textwidth}{| >{\raggedright\arraybackslash}p{3.9cm} | >{\raggedright\arraybackslash\small}p{5.3cm} | X |} \hline
\rowcolor{gray!33} Object Field & Value Type & Description \\ \hline
"system\allowbreak{}Name" & \pref{String} & Name of the system, PascalCase (e.g. \texttt{TemperatureSensor}). \\ \hline
"raw\allowbreak{}Key" & \pref{String} & The registered encryption key. \\ \hline
"algorithm" & \pref{String} & Encryption algorithm of the key, e.g. \texttt{AES/\allowbreak{}ECB/\allowbreak{}PKCS5Padding} or \texttt{AES/\allowbreak{}CBC/\allowbreak{}PKCS5Padding}. \\ \hline
"key\allowbreak{}Additive" & \pref{String} & Initialization vector (CBC mode), Base64; empty for ECB. \\ \hline
"created\allowbreak{}At" & \pref{DateTime} & UTC time of creation. \\ \hline
\end{tabularx}



\msubsection{struct}{AuthorizationTokenGenerationMgmtListRequest}

Tokens to generate on behalf of consumers.


\noindent\begin{tabularx}{\textwidth}{| >{\raggedright\arraybackslash}p{3.9cm} | >{\raggedright\arraybackslash\small}p{5.3cm} | X |} \hline
\rowcolor{gray!33} Object Field & Value Type & Description \\ \hline
"list" & \pref{Array}$<$\allowbreak{}\mref[Authorization\allowbreak{}Token\allowbreak{}Generation\allowbreak{}Mgmt\allowbreak{}Request]{AuthorizationTokenGenerationMgmtRequest}$>$ & The list of request items. \\ \hline
\end{tabularx}



\msubsection{struct}{AuthorizationTokenGenerationMgmtRequest}

One token to generate on behalf of a consumer.


\noindent\begin{tabularx}{\textwidth}{| >{\raggedright\arraybackslash}p{3.9cm} | >{\raggedright\arraybackslash\small}p{5.3cm} | X |} \hline
\rowcolor{gray!33} Object Field & Value Type & Description \\ \hline
"token\allowbreak{}Variant" & \pref{String} & Requested token variant, as interface security policy (e.g. \texttt{TIME\_\allowbreak{}LIMITED\_\allowbreak{}TOKEN\_\allowbreak{}AUTH}, \texttt{RSA\_\allowbreak{}SHA512\_\allowbreak{}JSON\_\allowbreak{}WEB\_\allowbreak{}TOKEN\_\allowbreak{}AUTH}). \\ \hline
"target\allowbreak{}Type" & \pref{String} & Target type: \texttt{SERVICE\_\allowbreak{}DEF} or \texttt{EVENT\_\allowbreak{}TYPE}. \\ \hline
"consumer\allowbreak{}Cloud" & \pref{String} & Cloud identifier of the consumer system. \\ \hline
"consumer" & \pref{String} & Name of the consumer system. \\ \hline
"provider" & \pref{String} & Name of the provider system. \\ \hline
"target" & \pref{String} & Target of the policy or token: a service definition or event type. \\ \hline
"scope" & \pref{String} & Operation the token or policy is limited to; empty for all operations. \\ \hline
"expires\allowbreak{}At" & \pref{DateTime} & Requested expiry; default \texttt{token.\allowbreak{}time.\allowbreak{}limit} from now. \\ \hline
"usage\allowbreak{}Limit" & \pref{Number} & Requested uses; default \texttt{simple.\allowbreak{}token.\allowbreak{}usage.\allowbreak{}limit}. \\ \hline
\end{tabularx}



\msubsection{struct}{AuthorizationTokenMgmtListResponse}

List of generated tokens.


\noindent\begin{tabularx}{\textwidth}{| >{\raggedright\arraybackslash}p{3.9cm} | >{\raggedright\arraybackslash\small}p{5.3cm} | X |} \hline
\rowcolor{gray!33} Object Field & Value Type & Description \\ \hline
"entries" & \pref{Array}$<$\allowbreak{}\mref[Authorization\allowbreak{}Token\allowbreak{}Response]{AuthorizationTokenResponse}$>$ & Matching records of the current page. \\ \hline
"count" & \pref{Number} & Total number of matching records (all pages). \\ \hline
\end{tabularx}



\msubsection{struct}{AuthorizationTokenQueryRequest}

Filter for generated tokens.


\noindent\begin{tabularx}{\textwidth}{| >{\raggedright\arraybackslash}p{3.9cm} | >{\raggedright\arraybackslash\small}p{5.3cm} | X |} \hline
\rowcolor{gray!33} Object Field & Value Type & Description \\ \hline
"pagination" & \mref{Page} & Paging and sorting of the result. \\ \hline
"requester" & \pref{String} & Name of the system that requested the token. \\ \hline
"token\allowbreak{}Type" & \pref{String} & Token type to match. \\ \hline
"consumer\allowbreak{}Cloud" & \pref{String} & Cloud identifier of the consumer system. \\ \hline
"consumer" & \pref{String} & Name of the consumer system. \\ \hline
"provider" & \pref{String} & Name of the provider system. \\ \hline
"target\allowbreak{}Type" & \pref{String} & Target type: \texttt{SERVICE\_\allowbreak{}DEF} or \texttt{EVENT\_\allowbreak{}TYPE}. \\ \hline
"target" & \pref{String} & Target of the policy or token: a service definition or event type. \\ \hline
\end{tabularx}



\msubsection{struct}{AuthorizationTokenResponse}

A generated token with its bookkeeping data.


\noindent\begin{tabularx}{\textwidth}{| >{\raggedright\arraybackslash}p{3.9cm} | >{\raggedright\arraybackslash\small}p{5.3cm} | X |} \hline
\rowcolor{gray!33} Object Field & Value Type & Description \\ \hline
"token\allowbreak{}Type" & \pref{AuthorizationTokenType} & Token type: \texttt{TIME\_\allowbreak{}LIMITED\_\allowbreak{}TOKEN}, \texttt{USAGE\_\allowbreak{}LIMITED\_\allowbreak{}TOKEN}, \texttt{SELF\_\allowbreak{}CONTAINED\_\allowbreak{}TOKEN} or \texttt{TRANSLATION\_\allowbreak{}BRIDGE\_\allowbreak{}TOKEN}. \\ \hline
"variant" & \pref{String} & Token variant (interface security policy). \\ \hline
"token" & \pref{String} & The token, for simple tokens only in hashed form. \\ \hline
"token\allowbreak{}Reference" & \pref{String} & Reference of the token, used to revoke it. \\ \hline
"requester" & \pref{String} & Name of the system that requested the token. \\ \hline
"consumer\allowbreak{}Cloud" & \pref{String} & Cloud identifier of the consumer system. \\ \hline
"consumer" & \pref{String} & Name of the consumer system. \\ \hline
"provider" & \pref{String} & Name of the provider system. \\ \hline
"target\allowbreak{}Type" & \pref{AuthorizationTargetType} & Target type: \texttt{SERVICE\_\allowbreak{}DEF} or \texttt{EVENT\_\allowbreak{}TYPE}. \\ \hline
"target" & \pref{String} & Target of the policy or token: a service definition or event type. \\ \hline
"scope" & \pref{String} & Operation the token or policy is limited to; empty for all operations. \\ \hline
"created\allowbreak{}At" & \pref{DateTime} & UTC time of creation. \\ \hline
"usage\allowbreak{}Limit" & \pref{Number} & Number of allowed uses (usage-limited tokens). \\ \hline
"usage\allowbreak{}Left" & \pref{Number} & Remaining uses of a usage-limited token. \\ \hline
"expires\allowbreak{}At" & \pref{DateTime} & UTC expiry time; no value means no expiry. \\ \hline
\end{tabularx}



\msubsection{struct}{ErrorMessage}




\noindent\begin{tabularx}{\textwidth}{| >{\raggedright\arraybackslash}p{3.9cm} | >{\raggedright\arraybackslash\small}p{5.3cm} | X |} \hline
\rowcolor{gray!33} Object Field & Value Type & Description \\ \hline
"error\allowbreak{}Message" & \pref{String} &  \\ \hline
"error\allowbreak{}Code" & \pref{Number} &  \\ \hline
"exception\allowbreak{}Type" & \pref{ExceptionType} &  \\ \hline
"origin" & \pref{String} &  \\ \hline
\end{tabularx}



\msubsection{struct}{Page}

Paging and sorting; all fields optional.


\noindent\begin{tabularx}{\textwidth}{| >{\raggedright\arraybackslash}p{3.9cm} | >{\raggedright\arraybackslash\small}p{5.3cm} | X |} \hline
\rowcolor{gray!33} Object Field & Value Type & Description \\ \hline
"page" & \pref{Number} & Zero-based page number. \\ \hline
"size" & \pref{Number} & Number of records per page (at most \texttt{max.\allowbreak{}page.\allowbreak{}size}, default 1000). \\ \hline
"direction" & \pref{String} & Sort direction, \texttt{ASC} or \texttt{DESC}. \\ \hline
"sort\allowbreak{}Field" & \pref{String} & Field to sort by. \\ \hline
\end{tabularx}



\subsection{Primitives}
\label{sec:model:primitives}

As all messages are encoded using the JSON format \cite{rfc8259}, the following primitive constructs, part of that standard, become available.
Note that the official standard is defined in terms of parsing rules, while this list only concerns syntactic information.
Furthermore, the \pref{Object} and \pref{Array} types are given optional generic type parameters, which are used in this document to signify when pair values or elements are expected to conform to certain types.

\noindent\begin{tabularx}{\textwidth}{| p{4.6cm} | X |} \hline
\rowcolor{gray!33} Type & Description \\ \hline
\pdef{Array} & A JSON Array \cite{rfc8259}; \pref{Array}$<$A$>$ has elements of type A. \\ \hline
\pdef{AuthorizationTargetType} & One of \texttt{SERVICE\_\allowbreak{}DEF}, \texttt{EVENT\_\allowbreak{}TYPE} (JSON String). \\ \hline
\pdef{AuthorizationTokenType} & One of \texttt{TIME\_\allowbreak{}LIMITED\_\allowbreak{}TOKEN}, \texttt{USAGE\_\allowbreak{}LIMITED\_\allowbreak{}TOKEN}, \texttt{SELF\_\allowbreak{}CONTAINED\_\allowbreak{}TOKEN}, \texttt{TRANSLATION\_\allowbreak{}BRIDGE\_\allowbreak{}TOKEN} (JSON String). \\ \hline
\pdef{DateTime} & JSON String with an ISO 8601 UTC timestamp, e.g. \texttt{2026-10-22T12:00:00Z}. \\ \hline
\pdef{ExceptionType} & One of \texttt{ARROWHEAD}, \texttt{INVALID\_\allowbreak{}PARAMETER}, \texttt{AUTH}, \texttt{FORBIDDEN}, \texttt{DATA\_\allowbreak{}NOT\_\allowbreak{}FOUND}, \texttt{TIMEOUT}, \texttt{LOCKED}, \texttt{INTERNAL\_\allowbreak{}SERVER\_\allowbreak{}ERROR}, \texttt{EXTERNAL\_\allowbreak{}SERVER\_\allowbreak{}ERROR} (JSON String). \\ \hline
\pdef{Number} & JSON Number, an integer. \\ \hline
\pdef{Object} & JSON Object; \pref{Object}$<$A$>$ has string keys and values of type A. \\ \hline
\pdef{String} & JSON String. \\ \hline
\end{tabularx}

\newpage

\bibliographystyle{IEEEtran}
\bibliography{bibliography}

\newpage

\section{Revision History}
\subsection{Amendments}

\noindent\begin{tabularx}{\textwidth}{| p{1cm} | p{3cm} | p{2cm} | X | p{4cm} |} \hline
\rowcolor{gray!33} No. & Date & Version & Subject of Amendments & Author \\ \hline

1 & 2026-10-22 & \arrowversion & Document created from the Arrowhead LaTeX templates for the 5.2.2 Java Spring implementation, including the local extensions of this workspace. & Jerker Delsing \\ \hline
\end{tabularx}

\subsection{Quality Assurance}

\noindent\begin{tabularx}{\textwidth}{| p{1cm} | p{3cm} | p{2cm} | X |} \hline
\rowcolor{gray!33} No. & Date & Version & Approved by \\ \hline

1 & & \arrowversion  &  \\ \hline

\end{tabularx}

\end{document}
